% =====================================================================
%  PARTE II. Reporte de finalización   (proceso TC de ISO/IEC/IEEE 29119-2)
% =====================================================================

\begin{notanormativa}[Correspondencia con el modelo de procesos]
Esta parte desarrolla el proceso de \textbf{finalización de prueba (TC1–TC4)} de \normap{2}:
archivar los activos reutilizables (TC1), limpiar y liberar el entorno (TC2), identificar
lecciones aprendidas (TC3) y reportar la finalización (TC4). Es, según el documento de procesos
consultado, el proceso más frecuentemente omitido en la práctica, pese a ser el que sostiene la
mejora continua.
\end{notanormativa}

\begin{notanormativa}[Seis preguntas distintas en el cierre]
El cierre responde a seis preguntas que conviene no mezclar, y que aquí ocupan apartados
distintos:
\begin{enumerate}
  \item qué actividades se completaron (\ref{sec:resumen-final});
  \item qué resultados se obtuvieron (\ref{sec:resultados});
  \item si se cumplieron los criterios de salida (\ref{sec:evaluacion-salida});
  \item qué incidentes quedan abiertos (\ref{sec:incidentes-abiertos});
  \item qué riesgos residuales persisten (\ref{sec:residuales});
  \item la decisión de aceptación y quién la toma (\ref{sec:aceptacion}).
\end{enumerate}
Haber ejecutado el 100\,\% de los casos —aun si muchos fallaron— no demuestra que el producto
sea aceptable.
\end{notanormativa}

\begin{observacion}[«Evaluación de la finalización» basada en cobertura]
\obsproblema{La plantilla anterior resolvía el cierre con un único apartado que determinaba «el
nivel de garantía de calidad logrado con base a la cobertura de las pruebas realizadas», sin
distinguir actividades completadas, resultados, incidentes abiertos, riesgo residual ni
autoridad de aceptación.}
\obscambio{El cierre se desglosa en los seis apartados enumerados arriba, con un cuadro de
decisión de aceptación que separa la recomendación técnica del equipo de la decisión de quien
asume el riesgo residual.}
\obsfuente{OBS §4.6}
\end{observacion}

% ---------------------------------------------------------------------
\section{Resumen de las pruebas realizadas}
\label{sec:resumen-final}

\begin{instruccion}
Compare las fechas planificadas con las realmente ejecutadas y presente el cronograma con
tiempos reales. Indique las incidencias que afectaron el progreso, sus causas y cómo afectaron
la ejecución del plan. Señale además qué actividades del Cuadro~\ref{tab:actividades} se
completaron, cuáles se completaron parcialmente y cuáles no se realizaron, con su motivo.
\end{instruccion}

\begin{table}[H]
\centering
\caption{Ejecución real frente a lo planificado.}
\label{tab:real-vs-plan}
\begin{tabularx}{\textwidth}{|C{1.4cm}|L{3.2cm}|C{2.0cm}|C{2.0cm}|C{1.9cm}|Y|}
  \hline
  \filaenc \textbf{Clave} & \textbf{Actividad} & \textbf{Fin planificado} & \textbf{Fin real} & \textbf{Estado} & \textbf{Causa de la desviación} \\ \hline
  \campo{ACT-01} & \campo{Actividad} & \campo{dd/mm} & \campo{dd/mm} & \campo{Completada} & \campo{---} \\ \hline
  & & & & & \\ \hline
\end{tabularx}
\notatabla{Estado: \textit{Completada} · \textit{Parcial} · \textit{No realizada} · \textit{Cancelada}.}
\end{table}

El Cuadro~\ref{tab:real-vs-plan} documenta qué se hizo realmente y por qué se desvió del plan.
Responde únicamente a la pregunta de si las actividades se completaron; no dice todavía nada
sobre la calidad del producto.

% ---------------------------------------------------------------------
\section{Resultados obtenidos}
\label{sec:resultados}

\begin{instruccion}
Presente los resultados agregados de la ejecución, tomados del \ref{anx:ejecuciones} y del
Cuadro~\ref{tab:cobertura-base}. El cuadro está dividido en los tres bloques que el
apartado~\ref{sec:cobertura} mantiene separados —avance, cobertura y veredicto—; respete la
división al rellenarlo y no traslade cifras de un bloque a otro.

En particular, la cobertura se calcula sobre lo \emph{ejercitado}: un elemento probado por un
caso que falló cuenta en el bloque de cobertura \textbf{y} en el de resultados. No lo descuente
del primero.
\end{instruccion}

\begin{table}[H]
\centering
\caption{Resultados agregados de la ejecución.}
\label{tab:resultados}
\begin{tabularx}{\textwidth}{|L{5.2cm}|C{2.6cm}|Y|}
  \hline
  \filaenc \textbf{Indicador} & \textbf{Valor} & \textbf{Comentario} \\ \hline
  \multicolumn{3}{|l|}{\cellcolor{uvgrisc}\textbf{Avance} — ¿cuánto del trabajo previsto se realizó?} \\ \hline
  Casos de prueba especificados        & \campo{n} & \campo{} \\ \hline
  Ejecuciones realizadas               & \campo{n} & \campo{Incluye repeticiones y regresiones} \\ \hline
  Casos bloqueados o no ejecutados     & \campo{n} & \campo{Motivo} \\ \hline
  Avance de ejecución (\%)             & \campo{\%} & Casos ejecutados ÷ casos planificados \\ \hline
  \multicolumn{3}{|l|}{\cellcolor{uvgrisc}\textbf{Cobertura} — ¿qué elementos quedaron ejercitados?} \\ \hline
  \campo{COB-01} — cobertura alcanzada & \campo{\%} & Elementos ejercitados, \emph{con independencia del veredicto}; fuente: Cuadro~\ref{tab:cobertura-base} \\ \hline
  \campo{COB-02} — cobertura alcanzada & \campo{\%} & \campo{Según su definición en el Cuadro~\ref{tab:cobertura}} \\ \hline
  Elementos del alcance sin ejercitar  & \campo{n} & \campo{Brechas de cobertura; deben reaparecer en \ref{sec:residuales}} \\ \hline
  \multicolumn{3}{|l|}{\cellcolor{uvgrisc}\textbf{Resultados} — de lo ejercitado, ¿qué se comportó como se esperaba?} \\ \hline
  Veredicto \emph{aprobado}            & \campo{n} & \campo{} \\ \hline
  Veredicto \emph{no aprobado}         & \campo{n} & \campo{Origina los incidentes del \ref{anx:incidentes}} \\ \hline
  Veredicto \emph{no concluyente}      & \campo{n} & \campo{} \\ \hline
\end{tabularx}
\end{table}

El Cuadro~\ref{tab:resultados} separa en tres bloques el avance, la cobertura y los veredictos.
Leídos juntos
responden a la pregunta que interesa al cierre: cuánto se probó, sobre qué y con qué resultado.
Ninguno de los tres bloques basta por sí solo para sostener la decisión de aceptación del
apartado~\ref{sec:aceptacion}.

% ---------------------------------------------------------------------
\section{Evaluación frente a los criterios de salida}
\label{sec:evaluacion-salida}

\begin{instruccion}
Evalúe uno a uno los criterios de finalización definidos en el
Cuadro~\ref{tab:criterios-salida}, indicando el valor alcanzado y si el criterio se cumple. No
reformule aquí los criterios: si alguno resultó inaplicable, declárelo como desviación en el
Cuadro~\ref{tab:desviaciones}.

Tome cada valor del bloque que le corresponde en el Cuadro~\ref{tab:resultados}: los criterios
de avance, de cobertura y de resultados se evalúan con cifras distintas. Un criterio de
cobertura cumplido al 100\,\% no compensa un criterio de resultados incumplido.
\end{instruccion}

\begin{instruccion}[Advertencia]
Si algún criterio de salida quedó redactado sobre «casos aprobados» pero clasificado como
criterio de \emph{cobertura}, corríjalo antes de evaluarlo: pertenece al tipo
\emph{Resultados}.
\end{instruccion}

\begin{table}[H]
\centering
\caption{Cumplimiento de los criterios de finalización.}
\label{tab:eval-salida}
\begin{tabularx}{\textwidth}{|C{1.5cm}|C{2.2cm}|Y|C{1.9cm}|C{1.9cm}|C{1.9cm}|}
  \hline
  \filaenc \textbf{Criterio} & \textbf{Tipo} & \textbf{Enunciado} & \textbf{Valor objetivo} & \textbf{Valor alcanzado} & \textbf{¿Se cumple?} \\ \hline
  \campo{CS-01} & \campo{Cobertura} & \campo{Enunciado} & \campo{Objetivo} & \campo{Real} & \campo{Sí/No} \\ \hline
  \campo{CS-02} & \campo{Avance} & \campo{Enunciado} & \campo{Objetivo} & \campo{Real} & \campo{Sí/No} \\ \hline
  \campo{CS-03} & \campo{Resultados} & \campo{Enunciado} & \campo{Objetivo} & \campo{Real} & \campo{Sí/No} \\ \hline
  \campo{CS-04} & \campo{Incidentes} & \campo{Enunciado} & \campo{Objetivo} & \campo{Real} & \campo{Sí/No} \\ \hline
  & & & & & \\ \hline
\end{tabularx}
\notatabla{La columna «Tipo» debe coincidir con la del Cuadro~\ref{tab:criterios-salida}. Sirve para comprobar de un vistazo que el cierre no se sostiene sobre un solo tipo de criterio.}
\end{table}

El Cuadro~\ref{tab:eval-salida} establece si las \emph{actividades de prueba} pueden darse por
concluidas, y la columna «Tipo» deja ver sobre qué se apoya esa conclusión. Es un insumo de la
decisión de aceptación del apartado~\ref{sec:aceptacion}, pero no la sustituye: cumplir los
criterios de salida significa que se hizo lo previsto, no que el resultado sea aceptable.

% ---------------------------------------------------------------------
\section{Incidentes abiertos al cierre}
\label{sec:incidentes-abiertos}

\begin{instruccion}
Relacione los incidentes que permanecen sin cerrar, con su severidad, prioridad y estado. Estos
incidentes son el principal insumo de los riesgos residuales y de la decisión de aceptación:
omitirlos convierte el cierre en una formalidad.
\end{instruccion}

\begin{table}[H]
\centering
\caption{Incidentes abiertos o diferidos al cierre del esfuerzo de prueba.}
\label{tab:incidentes-abiertos}
\begin{tabularx}{\textwidth}{|C{1.5cm}|Y|C{1.8cm}|C{1.8cm}|C{1.9cm}|C{2.2cm}|}
  \hline
  \filaenc \textbf{Clave} & \textbf{Resumen} & \textbf{Severidad} & \textbf{Prioridad} & \textbf{Estado} & \textbf{Disposición acordada} \\ \hline
  \campo{INC-01} & \campo{Resumen} & \campo{Alta} & \campo{Normal} & \campo{Diferido} & \campo{Qué se hará después} \\ \hline
  & & & & & \\ \hline
\end{tabularx}
\end{table}

El Cuadro~\ref{tab:incidentes-abiertos} es el extracto de cierre del \ref{anx:incidentes}; cada
fila debería tener su correspondencia en el Cuadro~\ref{tab:riesgos-residuales}.

% ---------------------------------------------------------------------
\section{Métricas de prueba}
\label{sec:metricas-finales}

\begin{instruccion}
Presente el cálculo final de las métricas definidas en el Cuadro~\ref{tab:metricas}, junto con
su análisis y discusión. Interprete cada valor según lo acordado en su definición y señale
explícitamente lo que la métrica \emph{no} permite concluir.
\end{instruccion}

\begin{table}[H]
\centering
\caption{Valores finales de las métricas de prueba.}
\label{tab:metricas-finales}
\begin{tabularx}{\textwidth}{|C{1.5cm}|C{1.7cm}|C{2.0cm}|Y|}
  \hline
  \filaenc \textbf{Clave} & \textbf{Tipo} & \textbf{Valor final} & \textbf{Análisis e interpretación} \\ \hline
  \campo{MET-01} & \campo{Producto} & \campo{Valor} & \campo{Qué indica y qué no} \\ \hline
  & & & \\ \hline
\end{tabularx}
\end{table}

El Cuadro~\ref{tab:metricas-finales} cierra el ciclo abierto en el apartado~\ref{sec:metricas}:
las métricas se interpretan según el criterio acordado \emph{antes} de conocer los resultados.

% ---------------------------------------------------------------------
\section{Riesgos residuales}
\label{sec:residuales}

\begin{instruccion}
Analice qué riesgos pudieron resolverse o evitarse y cuáles persisten y podrían afectar
procesos posteriores. Incluya al menos tres fuentes: los riesgos de producto no mitigados
(Cuadro~\ref{tab:riesgos-producto}), las exclusiones del alcance
(Cuadro~\ref{tab:exclusiones}) y los incidentes abiertos
(Cuadro~\ref{tab:incidentes-abiertos}).
\end{instruccion}

\begin{table}[H]
\centering
\caption{Riesgos residuales al cierre del esfuerzo de prueba.}
\label{tab:riesgos-residuales}
\begin{tabularx}{\textwidth}{|C{1.5cm}|C{2.0cm}|Y|C{1.9cm}|C{2.6cm}|}
  \hline
  \filaenc \textbf{Clave} & \textbf{Origen} & \textbf{Riesgo que persiste y su posible consecuencia} & \textbf{Exposición} & \textbf{Recomendación} \\ \hline
  \campo{RR-01} & \campo{RP-04 / EX-01 / INC-02} & \campo{Descripción} & \campo{Alta} & \campo{Acción sugerida} \\ \hline
  & & & & \\ \hline
\end{tabularx}
\end{table}

El Cuadro~\ref{tab:riesgos-residuales} declara lo que el esfuerzo de prueba \emph{no} llegó a
descartar. Es la información que la autoridad de aceptación necesita para decidir con
conocimiento de causa.

% ---------------------------------------------------------------------
\section{Decisión de aceptación}
\label{sec:aceptacion}

\begin{notanormativa}[Finalizar no es aceptar]
Dar por terminadas las actividades de prueba (\ref{sec:evaluacion-salida}) es una decisión
sobre el \emph{proceso}. Aceptar el producto es una decisión sobre el \emph{riesgo residual} que
la organización está dispuesta a asumir. Son decisiones distintas, se toman con información
distinta y, con frecuencia, corresponden a personas distintas, por lo que el plan debe nombrar
una autoridad de aceptación.
\end{notanormativa}

\begin{instruccion}
Registre la decisión, quién la toma y sobre qué base. La decisión debe considerar los resultados
obtenidos, la cobertura relevante alcanzada, los incidentes abiertos con su severidad y los
riesgos residuales. Si el equipo de prueba sólo emite una recomendación y la decisión
corresponde a otra parte, dígalo así.
\end{instruccion}

\begin{table}[H]
\centering
\caption{Decisión de aceptación del resultado de la prueba.}
\label{tab:aceptacion}
\begin{tabularx}{\textwidth}{|E{4.6cm}|Y|}
  \hline
  Recomendación del equipo de prueba & \campo{Aceptar / aceptar con condiciones / no aceptar} \\ \hline
  Base de la recomendación & \campo{Resultados, cobertura, incidentes abiertos y riesgos residuales considerados} \\ \hline
  Condiciones propuestas & \campo{Qué debería ocurrir antes de usar el software} \\ \hline
  Autoridad de aceptación & \autoridadAceptacion \\ \hline
  Decisión tomada & \campo{Decisión} \\ \hline
  Fecha & \campo{dd/mm/aaaa} \\ \hline
\end{tabularx}
\end{table}

El Cuadro~\ref{tab:aceptacion} separa explícitamente la recomendación técnica del equipo de
prueba de la decisión de quien tiene autoridad para asumir el riesgo residual.

\begin{politicaacademica}
En el contexto del curso, la autoridad de aceptación es \autoridadAceptacion, definida en
\texttt{config/metadata.tex}. Si el ejercicio académico no contempla una decisión real de
aceptación, indíquelo así en el cuadro y limite el contenido a la recomendación del equipo: es
preferible declarar que la decisión no aplica antes que simularla.
\end{politicaacademica}

% ---------------------------------------------------------------------
\section{Entregables de la prueba}

\begin{instruccion}
Relacione los entregables efectivamente producidos, comprometidos en el
apartado~\ref{sec:entregables} (Cuadro~\ref{tab:entregables}), indicando su
versión final y dónde se encuentran. Los procedimientos y casos de prueba aplicados se integran
en los anexos de este documento.
\end{instruccion}

\begin{table}[H]
\centering
\caption{Entregables producidos y su ubicación final.}
\label{tab:entregables-finales}
\begin{tabularx}{\textwidth}{|L{4.6cm}|C{2.0cm}|Y|C{1.9cm}|}
  \hline
  \filaenc \textbf{Entregable} & \textbf{Versión final} & \textbf{Ubicación} & \textbf{Estado} \\ \hline
  \campo{Entregable} & \campo{v1.0} & \campo{Ubicación} & \campo{Entregado} \\ \hline
  & & & \\ \hline
\end{tabularx}
\end{table}

El Cuadro~\ref{tab:entregables-finales} permite comprobar que lo prometido en el
Cuadro~\ref{tab:entregables} se produjo efectivamente.

% ---------------------------------------------------------------------
\section{Archivo, limpieza y activos reutilizables}
\label{sec:archivo}

\begin{notanormativa}[Correspondencia con el modelo de procesos]
Este apartado desarrolla las actividades \textbf{TC1 «Archivar los activos»} y \textbf{TC2
«Limpiar el entorno»} de \normap{2}: no basta con identificar qué activos son reutilizables, hay
que decir dónde se conservan, quién responde por ellos, qué ocurre con los datos sensibles y en
qué estado se deja el entorno.
\end{notanormativa}

\begin{observacion}[Archivo y limpieza sin custodia]
\obsproblema{La plantilla anterior pedía identificar los activos reutilizables, pero no su
ubicación, su responsable ni su periodo de conservación, y no contemplaba la liberación del
entorno ni el tratamiento de los datos sensibles al cierre.}
\obscambio{Dos cuadros: custodia de los activos archivados (TC1) y limpieza, liberación de
recursos, revocación de accesos y eliminación segura de datos (TC2).}
\obsfuente{OBS §3}
\end{observacion}

\begin{instruccion}
Complete los dos cuadros siguientes. El primero registra qué se conserva y bajo la
responsabilidad de quién; el segundo, qué se desmonta, libera o elimina. Ambos deben ser
coherentes con el control de configuración del Cuadro~\ref{tab:configuracion}.
\end{instruccion}

\begin{table}[H]
\centering
\caption{Activos de prueba conservados y su custodia.}
\label{tab:archivo}
\begin{tabularx}{\textwidth}{|L{3.4cm}|C{1.9cm}|Y|C{2.1cm}|C{2.0cm}|}
  \hline
  \filaenc \textbf{Activo} & \textbf{¿Reutilizable?} & \textbf{Dónde se archiva} & \textbf{Responsable} & \textbf{Periodo de conservación} \\ \hline
  \campo{Casos de prueba} & \campo{Sí} & \campo{Repositorio} & \campo{Rol} & \campo{Plazo} \\ \hline
  \campo{Procedimientos} & \campo{Sí} & \campo{Repositorio} & \campo{Rol} & \campo{Plazo} \\ \hline
  \campo{Guiones automatizados} & \campo{Sí/No} & \campo{Repositorio} & \campo{Rol} & \campo{Plazo} \\ \hline
  \campo{Evidencias de ejecución} & \campo{---} & \campo{Repositorio} & \campo{Rol} & \campo{Plazo} \\ \hline
  \campo{Registros de incidentes} & \campo{---} & \campo{Repositorio} & \campo{Rol} & \campo{Plazo} \\ \hline
  \campo{Datos de prueba} & \campo{Sí/No} & \campo{Repositorio} & \campo{Rol} & \campo{Plazo} \\ \hline
\end{tabularx}
\end{table}

\begin{table}[H]
\centering
\caption{Limpieza del entorno y liberación de recursos.}
\label{tab:limpieza}
\begin{tabularx}{\textwidth}{|L{3.8cm}|Y|C{2.1cm}|C{2.0cm}|}
  \hline
  \filaenc \textbf{Acción} & \textbf{Qué se hace concretamente} & \textbf{Responsable} & \textbf{Fecha} \\ \hline
  Restaurar el entorno & \campo{Estado al que se devuelve} & \campo{Rol} & \campo{dd/mm} \\ \hline
  Liberar recursos y licencias & \campo{Qué se libera} & \campo{Rol} & \campo{dd/mm} \\ \hline
  Tratar los datos sensibles & \campo{Eliminación segura, anonimización o devolución} & \campo{Rol} & \campo{dd/mm} \\ \hline
  Retirar accesos y credenciales & \campo{Qué se revoca} & \campo{Rol} & \campo{dd/mm} \\ \hline
\end{tabularx}
\notatabla{Si algún conjunto de datos contenía información personal, la eliminación segura no es opcional: documente el procedimiento aplicado.}
\end{table}

Los Cuadros~\ref{tab:archivo} y~\ref{tab:limpieza} distinguen lo que se \emph{conserva} de lo
que se \emph{desmonta}; las mejoras que se deriven de este cierre se registran en el
apartado~\ref{sec:lecciones}. Sin el segundo, un entorno de prueba con datos reales puede sobrevivir
al proyecto sin dueño ni control.

% ---------------------------------------------------------------------
\section{Lecciones aprendidas}
\label{sec:lecciones}

\begin{instruccion}
Discuta los resultados y plantee las conclusiones principales sobre el producto y sobre el
proceso. Evalúe el propio proceso de prueba e identifique áreas de mejora. Liste
recomendaciones justificadas y trazables hasta los resultados que las sustentan.

Indique, para cada lección, a quién se comunica y qué debería cambiar como consecuencia. Según
el modelo de procesos (actividad TC3), las lecciones aprendidas retroalimentan la capa
organizacional, tal como muestra el bucle externo de la Figura~\ref{fig:procesos}; una lección
que no llega a nadie no cierra ese bucle.
\end{instruccion}

\begin{table}[H]
\centering
\caption{Lecciones aprendidas y mejoras propuestas.}
\label{tab:lecciones}
\begin{tabularx}{\textwidth}{|C{1.4cm}|C{1.8cm}|Y|C{2.6cm}|C{2.0cm}|}
  \hline
  \filaenc \textbf{Clave} & \textbf{Ámbito} & \textbf{Lección y evidencia que la sustenta} & \textbf{Mejora propuesta} & \textbf{Se comunica a} \\ \hline
  \campo{LA-01} & \campo{Proceso} & \campo{Qué se aprendió y con qué evidencia} & \campo{Qué debería cambiar} & \campo{Destinatario} \\ \hline
  \campo{LA-02} & \campo{Producto} & \campo{Lección} & \campo{Mejora} & \campo{Destinatario} \\ \hline
  \campo{LA-03} & \campo{Proyecto} & \campo{Lección} & \campo{Mejora} & \campo{Destinatario} \\ \hline
  & & & & \\ \hline
\end{tabularx}
\notatabla{Ámbito: \textit{producto} (qué aprendimos sobre el software) · \textit{proceso} (sobre cómo probamos) · \textit{proyecto} (sobre cómo nos organizamos).}
\end{table}

El Cuadro~\ref{tab:lecciones} exige que cada lección esté sustentada en evidencia registrada y
tenga un destinatario concreto; así el cierre produce información utilizable en lugar de una
lista de buenas intenciones.
